% Part: turing-machines
% Chapter: undecidability
% Section: enumerating-tms

\documentclass[../../../include/open-logic-section]{subfiles}

\begin{document}

\olfileid{tur}{und}{enu}
\olsection{টুরিং যন্ত্রের তালিকায়ন}

\begin{explain}
সব টুরিং যন্ত্রের সেটটি যে !!{enumerable}, তা আমরা দেখাতে পারি। প্রতিটি
টুরিং যন্ত্রকে সসীমভাবে বর্ণনা করা যায় বলেই এই ফল পাওয়া যায়। দশার সেট ও
ফিতার বর্ণমালা সসীম সেট। অবস্থান্তর অপেক্ষকটি $Q \times \Sigma$ থেকে
$Q \times \Sigma \times \{\TMleft, \TMright, \TMstay\}$-এ একটি আংশিক
অপেক্ষক; তাই যে সসীমসংখ্যক আর্গুমেন্ট-যুগলে এটি সংজ্ঞায়িত, তাদের জন্য এর
মানগুলি তালিকায় লিখেও একে নির্দিষ্ট করা যায়।

এ পর্যন্ত কথাটি সত্য, কিন্তু এখানে একটি সূক্ষ্ম পার্থক্য আছে। টুরিং
যন্ত্রের সংজ্ঞায় দশার সেট ও ফিতার বর্ণমালায় কী ধরনের !!{element}s থাকতে
পারে, তার উপর কোনো বিধিনিষেধ দেওয়া হয়নি। তাই, যেমন, প্রতিটি বাস্তব
সংখ্যার জন্য এমন একটি টুরিং যন্ত্র প্রযুক্তিগতভাবে আছে, যা ওই সংখ্যাটিকে
একটি দশা হিসেবে ব্যবহার করে। কিন্তু কোন বস্তুগুলি দশা ও বর্ণমালার সদস্য
হিসেবে কাজ করছে, তার উপর টুরিং যন্ত্রটির \emph{আচরণ} নির্ভর করে না।
\olref{fig:variants}-এর দুটি টুরিং যন্ত্র বিবেচনা করো।
\begin{figure}
\begin{center}
\begin{tikzpicture}[->,>=stealth',shorten >=1pt,auto,node distance=2.8cm,
                    semithick]
  \tikzstyle{every state}=[fill=none,draw=black,text=black]

  \node[initial,state]         (A)              {$q_0$};
  \node[state]         (B) [right of=A] {$q_1$};

  \path (A) edge [bend left] node {\TMtrans{\TMstroke}{\TMstroke}{\TMright}} (B)
        (B) edge [loop above] node {\TMtrans{\TMblank}{\TMblank}{\TMright}} (B)
            edge [bend left] node {\TMtrans{\TMstroke}{\TMstroke}{\TMright}} (A);
\end{tikzpicture}\\
\begin{tikzpicture}[->,>=stealth',shorten >=1pt,auto,node distance=2.8cm,
  semithick]
\tikzstyle{every state}=[fill=none,draw=black,text=black]

\node[initial,state]         (A)              {$s$};
\node[state]         (B) [right of=A] {$h$};

\path (A) edge [bend left] node {\TMtrans{A}{A}{\TMright}} (B)
(B) edge [loop above] node {\TMtrans{\TMblank}{\TMblank}{\TMright}} (B)
edge [bend left] node {\TMtrans{A}{A}{\TMright}} (A);
\end{tikzpicture}
\end{center}
\caption{\emph{জোড়} যন্ত্রের রূপভেদ}
\ollabel{fig:variants}
\end{figure}
এই দুটি চিত্র দুটি যন্ত্রের নির্দেশ দেয়: $M$-এর ফিতা-বর্ণমালা
$\Sigma = \{\TMendtape,\TMblank,\TMstroke\}$ এবং দশা-সেট
$\{q_0,q_1\}$; আর $M'$-এর বর্ণমালা $\Sigma' =
\{\TMendtape,\TMblank,A\}$ এবং দশা-সেট $\{s,h\}$। কিন্তু অন্য সব দিক
থেকে তাদের নির্দেশ একই: $n$ জোড় হলে এবং কেবল তখনই $M$ $n$টি
$\TMstroke$-এর প্রতীকক্রমে থামবে, আর $n$ জোড় হলে এবং কেবল তখনই $M'$
$n$টি~$A$-এর প্রতীকক্রমে থামবে। আমরা শুধু $\TMstroke$-এর নতুন নাম~$A$,
$q_0$-র নতুন নাম~$s$, এবং $q_1$-এর নতুন নাম~$h$ দিয়েছি। এই উদাহরণটি
সাধারণীকরণ করা যায়: প্রতীক ও দশার এমন নামবদল করে একটি থেকে অন্যটি পাওয়া
গেলে টুরিং যন্ত্রদুটিকে একই বলে ভাবতে পারি। আসলে টুরিং যন্ত্রের প্রতীক ও
দশাগুলিকে আমরা সরাসরি ধনাত্মক পূর্ণসংখ্যা হিসেবে ভাবতে পারি: $\sigma_0$-র
বদলে~$1$, $\sigma_1$-এর বদলে~$2$, ইত্যাদি; $\TMendtape$ হলো~$1$,
$\TMblank$ হলো~$2$, ইত্যাদি। এভাবে \emph{জোড়} যন্ত্রটি
\olref{fig:standard-even}-এর চিত্রের যন্ত্রে পরিণত হয়।
\begin{figure}
\[\begin{tikzpicture}[->,>=stealth',shorten >=1pt,auto,node distance=2.8cm,
  semithick]
\tikzstyle{every state}=[fill=none,draw=black,text=black]

\node[initial,state]         (A)              {$1$};
\node[state]         (B) [right of=A] {$2$};

\path (A) edge [bend left] node {\TMtrans{3}{3}{\TMright}} (B)
(B) edge [loop above] node {\TMtrans{2}{2}{\TMright}} (B)
edge [bend left] node {\TMtrans{3}{3}{\TMright}} (A);
\end{tikzpicture}
\]
\caption{একটি মানক \emph{জোড়} যন্ত্র}
\ollabel{fig:standard-even}
\end{figure}
যে টুরিং যন্ত্রের দশা ও প্রতীক ধনাত্মক পূর্ণসংখ্যা, তাকে \emph{মানক}
যন্ত্র বলতে পারি এবং এখন থেকে শুধু মানক যন্ত্রই বিবেচনা করতে পারি।
\footnote{“মানক যন্ত্র” পরিভাষাটি মানক নয়।}

আমরা দেখাতে চেয়েছিলাম যে টুরিং যন্ত্রের সেটটি !!{enumerable}; উপরের
বিবেচনা অনুসারে মানক টুরিং যন্ত্রের সেটটি !!{enumerable} দেখালেই যথেষ্ট।
ধরা যাক, একটি মানক টুরিং যন্ত্র $M = \tuple{Q, \Sigma, q_0, \delta}$
দেওয়া আছে। ধনাত্মক পূর্ণসংখ্যার একটি সসীম প্রতীকক্রম দিয়ে একে কীভাবে
বর্ণনা করব? প্রথমে দশার সংখ্যা, দশাগুলি, প্রতীকের সংখ্যা, প্রতীকগুলি এবং
আরম্ভিক দশা তালিকায় লিখব। (মনে রেখো, $M$ মানক যন্ত্র বলে এগুলি সবই
ধনাত্মক পূর্ণসংখ্যা।) এবার~$\delta$-র কী হবে? সম্ভাব্য আর্গুমেন্টের সেট,
অর্থাৎ $\tuple{q,\sigma}$ যুগলগুলির সেট, সসীম; কারণ $Q$ ও~$\Sigma$
সসীম। তাই~$\delta$-র তথ্যটি ঠিক সেই সব $5$-ক~$\tuple{q, \sigma, q',
\sigma', d}$-এর সসীম তালিকা, যাদের জন্য $\delta(q,\sigma) = \tuple{q',
\sigma', D}$; এখানে $d$ হলো দিক~$D$-এর সংকেত-সংখ্যা (ধরা যাক,
~$\TMleft$-এর জন্য $1$, $\TMright$-এর জন্য $2$, এবং~$\TMstay$-এর জন্য
$3$)।

এভাবে প্রতিটি মানক টুরিং যন্ত্রকে ধনাত্মক পূর্ণসংখ্যার একটি সসীম তালিকা,
অর্থাৎ একটি অনুক্রম $s_M \in (\PosInt)^*$ দিয়ে বর্ণনা করা যায়। যেমন,
মানক \emph{জোড়} যন্ত্রটির সংকেত হলো অনুক্রমটি
\[
2, \underbrace{1, 2}_Q, 3, \overbrace{1, 2, 3}^\Sigma, 1, \underbrace{1, 3, 2, 3, 2}_{\delta(1,3) = \tuple{2,3,R}}, 
\overbrace{2, 2, 2, 2, 2}^{\delta(2,2) = \tuple{2,2,R}},
\underbrace{2, 3, 1, 3, 2}_{\delta(2,3) = \tuple{1,3,R}}.
\]
\end{explain}

\begin{thm}
$\Nat$ থেকে~$\Nat$-এ এমন অপেক্ষক আছে যা টুরিং-গণনসাধ্য নয়।
\end{thm}

\begin{proof}
আমরা জানি, ধনাত্মক পূর্ণসংখ্যার সসীম অনুক্রমগুলির সেট~$(\PosInt)^*$
!!{enumerable} (\cref{sfr:siz:zigzag:prob:posint-star})। তাই মানক টুরিং
যন্ত্রের বর্ণনাগুলির সেটটি~$(\PosInt)^*$-এর একটি উপসেট হিসেবে নিজেও
তালিকায়নযোগ্য। $\Nat$ থেকে~$\Nat$-এ প্রতিটি টুরিং-গণনসাধ্য অপেক্ষক কোনো
না কোনো (আসলে বহু) টুরিং যন্ত্রে গণনা করা হয়। তার দশা ও প্রতীকগুলির নতুন
নাম ধনাত্মক পূর্ণসংখ্যা দিয়ে—বিশেষত $\TMendtape$-কে~$1$,
$\TMblank$-কে~$2$, এবং $\TMstroke$-কে~$3$ নাম দিয়ে—দেখা যায় যে প্রতিটি
টুরিং-গণনসাধ্য অপেক্ষক কোনো মানক টুরিং যন্ত্রে গণনা করা হয়। এর অর্থ,
$\Nat$ থেকে~$\Nat$-এ সকল টুরিং-গণনসাধ্য অপেক্ষকের সেটটিও তালিকায়নযোগ্য।

অন্যদিকে, $\Nat$ থেকে~$\Nat$-এ সকল অপেক্ষকের সেটটি !!{enumerable} নয়
(\cref{sfr:siz:red:prob:nat-nat})। সব অপেক্ষকই কোনো টুরিং যন্ত্রে গণনা
করা গেলে যেসব টুরিং যন্ত্র তাদের গণনা করে, সেগুলির বর্ণনা তালিকায় লিখে
সব অপেক্ষকের সেটটিকেও তালিকায়িত করতে পারতাম। তাই কিছু অপেক্ষক
টুরিং-গণনসাধ্য নয়। 
\end{proof}

\begin{prob}
  দশা ও বর্ণমালার প্রতীকগুলি স্পষ্টভাবে তালিকায় না-লিখেও টুরিং যন্ত্র
  বর্ণনা করার কোনো উপায় কি ভাবতে পারো? “মানক” যন্ত্রের নিজস্ব ধারণা
  সংজ্ঞায়িত করতে পারো; তবে তোমার নতুন অর্থে প্রতিটি টুরিং যন্ত্রকে কেন
  কোনো “মানক” যন্ত্র অনুকরণ করতে পারে, সে বিষয়ে কিছু বলো।
\end{prob}
% BN-SRC-175 TARGET-CORRECTION-BEGIN
\par\smallskip\noindent\textnormal{\small\textbf{উৎস-সংশোধন (BN-SRC-175):} হিমায়িত ইংরেজি অনুশীলনটি বলে প্রতিটি টুরিং যন্ত্রকে একটি মানক যন্ত্র “compute” করতে পারে, যদিও যন্ত্র অন্য যন্ত্রকে গণনা করে না; আগের নামবদল যুক্তি অনুসারে মানক যন্ত্রটি তার আচরণ অনুকরণ করে। বাংলা পাঠে সেই সম্পর্কটিই দেওয়া হয়েছে।} % BN-SRC-175 TARGET-CORRECTION-END

\end{document}
